\chapter{117}\label{section-117}

Der bissige Wind veranlasste Louis den Jackenkragen hochzuschlagen.
Nicht nur das Wetter liess ihn frösteln.

Die Bar lag am Ende eines geraden Wegstücks. Er wusste, sie war bereits
geschlossen. Saisonende ab November. Sie wirkte aus der Ferne
entsprechend ausgestorben.

Sie hatten sich draussen verabredet. Louis hatte sich mit Valentina auf
einen gemeinsamen Spaziergang geeinigt, hatte er Joh erklärt.

Da kam sie. Hatte sie sich die Haare kurz geschnitten? Louis` Körper
spannte sich an. Mit zusammengekniffenen Augen fokussierte er die Person
in der Ferne. Nein, das war nicht Valentina. Je näher er kam, desto
klarer konnte er eine Frau erkennen. Oder doch eher ein junges Mädchen.
Knapp erwachsen, schätzte er, als sie nah genug war. Sie schaute ihm
direkt ins Gesicht. Selbst als Louis vor ihr stand, hielt sie den Blick.
An Selbstsicherheit schien es ihr nicht zu mangeln.

Sie lächelte.

«Louis?»

«Ja,» er zog die Augenbrauen hoch.

Seine Verwirrung schien das Mädchen zu amüsieren.

«Hi,» sie kam ihm einen Schritt entgegen und reichte ihm ihre schmale
Hand, «ich bin Isabella. Valentina schickt mich. Sie ist krank.»

Louis konnte geradezu fühlen, wie er in sich zusammenfiel. Er wusste
nicht, worüber er sich mehr ärgerte. Ob über Valentinas selbstsichere
Stellvertreterin oder über die Tatsache, dass Valentina ihn versetzte.

«Valentina ist krank, wirklich?»

«Ja, ja, Migräne, eine starke Migräne. Sie lässt dich grüssen.»

Sie begann ungeduldig in ihrer Umhängetasche zu wühlen.

Louis fühlte Zorn aufsteigen. Sein Plan schien dahin. Ein Schatten legte
sich über sein Gesicht.

«Plötzlich krank? Gestern klang sie am Telefon noch ganz schön vital.
Wer bist du?»

Mit einer resoluten Bewegung hatte sie es endlich geschafft. Sie zog
einen Briefumschlag aus der Tasche und reichte ihn Louis.

«Ich bin Vals Kusine. Sie hat mir viel von dir erzählt.»

Sie schmunzelte geheimnisvoll.

«Ja, ihre Migräne. Die kommt plötzlich. Ich glaube, immer mal wieder,
wenn sie durcheinander ist.»

Jetzt zwinkerte sie ihm zu. Rätselhaft. Es wurde immer besser.

Wurde er gerade für blöd verkauft?

«Warum schickt sie \emph{dich?} Warum rief sie mich nicht einfach an?»

Die Worte waren bissiger über Louis` Lippen gekommen, als ihm recht war.
Die Kleine konnte ja nichts dafür.

Sie zuckte erst nur mit den Schultern, setzte aber plötzlich einen
andern Blick auf. War sie beleidigt?

Sie schien sich beherrschen zu müssen.

«Woher soll ich das wissen?» zischte sie, «frag sie doch selber?»

Sie wandte sich entschieden um und ging davon. Bevor sie ihr Fahrrad
erreicht hatte, drehte sie sich noch einmal um.

«Dich habe ich mir echt cooler vorgestellt. Und bitte auch,» rief sie
trotzig hinter sich, «sehr gern geschehen, ciao.»

Sie stieg auf ihr Rad und liess Louis mit dem Brief in der Hand zurück.
